\documentclass[a4paper,12pt]{article}
\usepackage[utf8]{inputenc}
\usepackage[italian]{babel}
\usepackage{listings}
\usepackage{xcolor}
\usepackage{float}
\usepackage{placeins}

\lstset{
  language=C,
  backgroundcolor=\color{lightgray!20},
  frame=single,
  breaklines=true,
  numbers=left,
  numberstyle=\tiny,
  stepnumber=1,
  numbersep=5pt,
  captionpos=b,
  basicstyle=\ttfamily\small,
  keywordstyle=\color{blue}\bfseries,
  commentstyle=\color{green!60!black},
  stringstyle=\color{red},
  showstringspaces=false,
  float=H
}

\title{LabReSiD26 - Hands-On 2}
\author{Davide Ferrara - 518629}
\date{}

\begin{document}

\maketitle

% ---------------------------------------------------------------------------
\section{Esercizio 1: Prova ad immaginare il funzionamento su rete reale}
% ---------------------------------------------------------------------------

\subsection{Dove si trova esattamente il problema della frammentazione nel client e nel server?}
Il problema non è dato dal codice in sé ma dal fatto che tra Client e Server non è stabilito
alcun protocollo di comunicazione. TCP è un protocollo che garantisce l'invio ordinato di un
flusso di byte (stream-oriented), non esiste il concetto di messaggio in TCP.

Dal manuale di TCP (\texttt{man 7 tcp}):

\begin{quote}
\textit{TCP guarantees that the data arrives in order and retransmits lost packets.
It generates and checks a per-packet checksum to catch transmission errors.
TCP does not preserve record boundaries.}
\end{quote}

\subsection{Perché avviene in trasmissione?}

TCP può spezzare i dati in più segmenti per via del MSS (Maximum Segment Size),
negoziato durante il 3-way handshake. Inoltre il Nagle's algorithm può raggruppare
più \texttt{send()} in un unico segmento per ridurre l'overhead di rete.

\subsection{Perché avviene in ricezione?}

\texttt{recv()} restituisce quanti byte sono disponibili nel receive buffer del kernel
in quel momento, senza aspettare un messaggio completo. Possono arrivare meno byte
del previsto (segmento frammentato) o più byte del previsto (due \texttt{send()} fusi
in un unico segmento).

% ---------------------------------------------------------------------------
\section{Esercizio 2: Scrivi \texttt{recv\_all()}}
% ---------------------------------------------------------------------------

\begin{lstlisting}[caption={recv\_all() -- implementazione}]
ssize_t recv_all(int sockfd, void *buf, size_t n) {
  ssize_t r = 0;

  /* cicla finche' non ha letto esattamente n byte */
  while ((size_t)r < n) {
    /* aritmetica di puntatori: (char*)buf + r sposta il cursore
       sui byte gia' ricevuti; n-r e' quanto manca ancora */
    ssize_t received = recv(sockfd, (char *)buf + r, n - r, 0);
    if (received == -1) { /* errore di sistema */
      perror("[ERROR] recv errore: ");
      return -1;
    }
    if (received == 0) { /* peer ha chiuso la connessione (FIN) */
      printf("[DEBUG] connessione chiusa!\n");
      return 0;
    }
    r += received; /* aggiorna il cursore */
  }

  return n; /* restituisce n byte letti con successo */
}
\end{lstlisting}

\begin{lstlisting}[caption={send\_all() -- implementazione}]
ssize_t send_all(int sockfd, void *buf, size_t n) {
  ssize_t s = 0;

  /* cicla finche' non ha inviato esattamente n byte */
  while ((size_t)s < n) {
    /* aritmetica di puntatori: (char*)buf + s sposta il cursore
       sui byte gia' inviati; n-s e' quanto manca ancora */
    ssize_t sent = send(sockfd, (char *)buf + s, n - s, 0);
    if (sent == -1) { /* errore di sistema */
      perror("[ERROR] send errore: ");
      return -1;
    }
    if (sent == 0) { /* connessione chiusa */
      printf("[DEBUG] connessione chiusa!\n");
      return 0;
    }
    s += sent; /* aggiorna il cursore */
  }

  return n; /* restituisce n byte inviati con successo */
}
\end{lstlisting}

L'implementazione di questa funzione non risolve il problema ugualmente perché i messaggi
sono di lunghezza variabile, serve inviare anche questa informazione al server dal client.

% ---------------------------------------------------------------------------
\section{Esercizio 3: Il problema della lunghezza}
% ---------------------------------------------------------------------------

\subsection{a) Il mittente deve mandare prima la lunghezza o prima il payload?}

Il mittente deve mandare prima la lunghezza, così il ricevitore legge i 4 byte di header,
conosce quanti byte aspettarsi e chiama \texttt{recv\_all()} con quella lunghezza esatta.

\subsection{b) Dimensione massima di un singolo messaggio con \texttt{uint32\_t}}

\texttt{uint32\_t} va da 0 a $2^{32}-1$, quindi la dimensione massima del payload è
$2^{32}-1$ byte, circa 4 GB.

\subsection{c) Perché è necessario usare \texttt{htonl()} sul campo lunghezza?}

Macchine diverse potrebbero usare endianness diverse (little-endian su x86, big-endian su
architetture a 64 bit). Il network byte order è big-endian per convenzione. Senza
\texttt{htonl()}, un client little-endian invia i byte in ordine inverso e il server
interpreta un numero completamente sbagliato.

% ---------------------------------------------------------------------------
\section{Esercizio 4: Implementa il protocollo}
% ---------------------------------------------------------------------------

La struct \texttt{message} usata da client e server:

\begin{lstlisting}[caption={Struttura message}]
typedef struct message {
  uint32_t len;
  char *payload;
} message;
\end{lstlisting}

Il mittente invia prima i 4 byte di lunghezza in network byte order (\texttt{htonl}),
poi il payload. Il ricevitore legge i 4 byte, converte con \texttt{ntohl}, alloca
esattamente \texttt{len+1} byte e chiama \texttt{recv\_all()} con quella lunghezza.

\begin{lstlisting}[caption={send\_message() -- client}]
int send_message(int sockfd, const message *msg) {
  uint32_t net_len = htonl(msg->len);
  if (send_all(sockfd, &net_len, sizeof(uint32_t)) < 0)
    return -1;
  if (send_all(sockfd, (void *)msg->payload, msg->len) < 0)
    return -1;
  return 0;
}
\end{lstlisting}

Il ricevente legge i primi 4 byte, la lunghezza, e li converte in host long,
alloca la memoria (len + 1) e termina la stringa.

\begin{lstlisting}[caption={recv\_message() -- server}]
int recv_message(int sockfd, message *msg) {
  if (recv_all(sockfd, &msg->len, sizeof(uint32_t)) <= 0)
    return -1;
  msg->len = ntohl(msg->len);
  msg->payload = (char *)malloc(sizeof(char) * msg->len + 1);
  if (recv_all(sockfd, msg->payload, msg->len) <= 0)
    return -1;
  msg->payload[msg->len] = '\0';
  return 0;
}
\end{lstlisting}

\textbf{Output} Con un messaggio piccolo \texttt{"Ciao col protocollo!"} (20 byte):

\begin{lstlisting}[caption={Output -- messaggio piccolo}, language={}]
[DEBUG] Lunghezza divina_commedia.txt: 64000 bytes
[CLIENT.c] Inviando: Ciao col protocollo!
[DEBUG] r=4
[DEBUG] r=20
\end{lstlisting}

I \texttt{[DEBUG] r=4} e \texttt{r=20} confermano che \texttt{send\_all()} trasmette prima
i 4 byte di header, poi i 20 byte di payload.

\textbf{Test con messaggio da 64 KB} Il client invia un testo della \textit{Divina Commedia}
(64\,000 byte) in un singolo messaggio. Il server riceve correttamente l'intero payload
grazie a \texttt{recv\_all()} che cicla fino a leggere esattamente \texttt{len} byte:

\begin{lstlisting}[caption={Output server -- payload 64 KB}, language={}]
Rispuosemi: "Non omo, omo gia' fui,
e li parenti miei furon lombardi,
mantoani per patria ambedui.
[...]
[DEBUG] connessione chiusa!
\end{lstlisting}

% ---------------------------------------------------------------------------
\section{Esercizio 5: Aggiungi il checksum}
% ---------------------------------------------------------------------------

Il checksum XOR viene calcolato su tutti i byte del payload prima dell'invio e verificato
dal server dopo la ricezione. Il formato del messaggio diventa:

\begin{center}
\texttt{| len (4 byte) | payload (len byte) | checksum (1 byte) |}
\end{center}

La struct \texttt{message} viene estesa con il campo \texttt{checksum}:

\begin{lstlisting}[caption={Struct message con checksum}]
typedef struct message {
  uint32_t len;
  char *payload;
  uint8_t checksum;
} message;
\end{lstlisting}

\newpage
Il checksum viene calcolato come lo XOR di tutti i caratteri del payload.
\begin{lstlisting}[caption={Calcolo checksum XOR}]
uint8_t checksum(char *payload) {
  uint8_t checksum = 0;
  for (size_t i = 0; i < strlen(payload); i++) {
    checksum ^= (uint8_t)payload[i];
  }
  return checksum;
}
\end{lstlisting}

Il client calcola il checksum in \texttt{create\_message()} e lo invia come ultimo byte
dopo il payload in \texttt{send\_message()}:

\begin{lstlisting}[caption={Invio checksum -- client}]
/* In create_message(): */
msg->checksum = checksum(msg->payload);

/* In send_message(): */
printf("[CLIENT.c] Checksum inviato: %d\n", msg->checksum);
if (send_all(sockfd, (void *)&msg->checksum, sizeof(uint8_t)) < 0)
  return -1;
\end{lstlisting}

Il server riceve il byte di checksum, calcola il proprio e confronta:

\begin{lstlisting}[caption={Verifica checksum -- server}]
if (recv_all(sockfd, &msg->checksum, sizeof(uint8_t)) <= 0)
  return -1;

uint8_t computed = checksum(msg->payload);
printf("[SERVER.c] Checksum ricevuto: %d | Checksum calcolato: %d\n",
       msg->checksum, computed);

if (!(checksum(msg->payload) && msg->checksum)) {
  printf("[DEBUG] Checksum fallito!\n");
  return -1;
}
printf("[DEBUG] Checksum corretto, il messaggio e' integro!\n");
\end{lstlisting}

Il server tronca il payload stampato a 64 caratteri per leggibilità
(stampare 64KB di divina commedia nel terminal sarebbe impossibile),
il checksum controlla comunque la validità del messaggio.

\begin{lstlisting}[caption={Troncamento output -- server}]
if (strlen(msg->payload) > 64) {
  msg->payload[61] = '.';
  msg->payload[62] = '.';
  msg->payload[63] = '.';
  msg->payload[64] = '\0';
}
\end{lstlisting}

\textbf{Output con file da 64\,000 byte} (\textit{Divina Commedia}):

\begin{lstlisting}[caption={Output -- checksum su 64 KB}, language={}]
[DEBUG] r=4
[DEBUG] r=64000
[DEBUG.c] len: 64000
[DEBUG] r=1
[SERVER.c] Checksum ricevuto: 54 | Checksum calcolato: 54
[DEBUG] Checksum corretto, il messaggio e' integro!
[SERVER.c] Messaggio: Nel mezzo del cammin di nostra vita
mi ritrovai per una selva...
\end{lstlisting}

\texttt{r=4} riceve il campo lunghezza, \texttt{r=64000} l'intero payload,
\texttt{r=1} il byte di checksum. I valori \texttt{54} coincidono: il messaggio
è arrivato integro.

% ---------------------------------------------------------------------------
\section{Domande di riflessione}
% ---------------------------------------------------------------------------

\subsection*{Q1. Perché il codice funzionava in locale ma può fallire su rete reale?}

In loopback il kernel consegna i dati in un unico chunk mentre nella rete reale TCP
può spezzarli.

\subsection*{Q2. Differenza tra \texttt{recv()} che ritorna 0 e \texttt{recv()} che ritorna -1}

\texttt{recv()} ritorna 0 se il client ha effettuato uno shutdown in ordine (FIN ricevuto).
\texttt{recv()} ritorna -1 nel caso di errori come una chiusura improvvisa della connessione.

\subsection*{Q3. Dimensione massima del payload con \texttt{uint16\_t} nel campo lunghezza}

\texttt{uint16\_t} corrisponde a un intero senza segno da 16 bit quindi max =
$2^{16}-1 = 65535$ byte di payload.

\subsection*{Q4. Il protocollo è adatto a trasferire file binari?}

Il protocollo in generale sì, ma non nel mio codice perché nel checksum uso
\texttt{strlen()}, che si ferma al primo byte \texttt{0x00}. I file binari contengono
null byte essendo una sequenza di 0 e 1, quindi il checksum sarebbe errato.
Anche la struttura \textbf{message} non sarebbe adatta per via del buffer di
tipo char*.

\end{document}
